\section{Architettura del Modello}
Il modello implementato in NetLogo discretizza lo spazio in una \texttt{world grid} di 41x41 patch, dove ciascuna patch può rappresentare uno spazio percorribile o un muro. Un pedone occupa esattamente una cella.

La topologia della rotonda è costruita intagliando un'area calpestabile: un anello circolare delimitato da un raggio interno e uno esterno, incrociato da due corridoi rettilinei ortogonali. La rottura della simmetria originata dall'isola centrale impone ai flussi perpendicolari di inserirsi obbligatoriamente nella circolazione tangenziale dell'anello.

La cinematica discreta del modello usa un \textit{Moore neighborhood} a 8 celle adiacenti più l'opzione di permanenza sulla cella corrente. Ciò produce 9 possibili scelte di transizione per ogni agente: 8 derivanti dalle direzioni cardinali/diagonali e 1 dalla possibilità di restare fermi (\texttt{stay}).

Questa scelta è coerente con il \textit{flood fill} dei campi statici e riduce l'anisotropia che emergerebbe da un vicinato di Von Neumann. Dal punto di vista fisico, il vicinato di Moore permette micro-correzioni diagonali e rende più naturale la circolazione nell'anello.

Questa è la ragione per cui il sistema può essere descritto come esempio di intelligenza distribuita passiva: l'ordine globale non viene imposto tramite controllo centralizzato, bensì tramite \textit{environmental biasing}.
